\documentclass[dvipdfmx, 11pt, a4paper]{jsarticle}

% --- 基本パッケージ ---
\usepackage[T1]{fontenc}
\usepackage{mathrsfs}
\usepackage{lmodern}
\usepackage{amsmath, amssymb, amsthm, mathtools}
\usepackage{bm}           % ベクトルなどの太字 \bm{}
\usepackage{graphicx}     % 画像の挿入
\usepackage{hyperref}     % リンク機能
\usepackage{footnotehyper}
\usepackage{ulem}         %斜線


\usepackage{url}

% --- 物理・数式用便利パッケージ ---
\usepackage{physics}      % \pdv, \dd, \bra, \ket などが使える神パッケージ
\usepackage{braket}
\usepackage{siunitx}      % 単位の記述 \si{\meter}

% --- 見た目の調整 (tcolorbox) ---
\usepackage[most]{tcolorbox}

% 定理・定義などの枠デザイン
\tcbset{
    common/.style={
        enhanced,
        colback=white,
        boxrule=0.5mm,
        sharp corners,
        fonttitle=\bfseries
    }
}

% 定理・公式（深いネイビー）
\newtcolorbox{definition}[1][定義]{
    common,
    colframe=blue!40!gray!75,    % グレーがかった青
    colback=blue!5!gray!5!white,
    title={#1}
}

% 補題・定義（濃いシアン・ティール）
\newtcolorbox{formula}[1][定理]{
    common,
    colframe=cyan!30!black,  % 黒に近い水色（青緑）
    colback=cyan!3!white,
    title={#1}
}


% 補足（モノトーンで邪魔しない）
\newtcolorbox{note}[1][Note]{
    common,
    colframe=black!40,        % 濃いグレー
    colback=gray!5!white,
    title={#1},
    borderline west={1mm}{0mm}{black!60}
}


% --- 証明環境の設定 ---
\renewcommand{\proofname}{\textbf{証明}}
\theoremstyle{definition}
\newtheorem{thm}{定理}
\newtheorem{lem}{補題}

% --- 余白の設定（少し広めに取る） ---
\usepackage[top=25mm, bottom=25mm, left=20mm, right=20mm]{geometry}

% --- 証明用の新しい環境定義 (左線スタイル) ---
\newtcolorbox{derivation}[1][導出]{
    empty,                  % 背景や枠を無効化
    breakable,              % ページまたぎOK
    title={#1},             % タイトル（デフォルトは「証明」）
    coltitle=black,         % タイトルの色
    fonttitle=\bfseries,    % タイトルのフォント
    left=5mm,               % 左の余白
    parbox=false,           % 通常の段落として扱う
    borderline west={1mm}{0pt}{gray!60}, % 左側に太さ1mmのグレー線
    before upper={\setlength{\parindent}{1em}}, % インデントを保持
}

\newtcolorbox{derivationBox}[1][\textbf{導出}]{
    enhanced,
    breakable,
    frame hidden,
    colback=gray!10!white,
    coltitle=black,
    fonttitle=\bfseries,
    title={#1},
    attach title to upper,       % タイトルを本文領域に埋め込む
    after title={\par\smallskip},% ★ 改行(\par)と、少しの隙間(\smallskip)を入れる
    left=3mm, right=3mm, top=3mm, bottom=3mm,
    parbox=false
}




\usepackage{fancyhdr}
\pagestyle{fancy}
\fancyhf{} % 一旦リセット
\lhead{\leftmark} % 左上にセクション名
\rhead{\thepage}  % 右上にページ番号
\cfoot{}          % 中央下は空にする
% jsarticleの場合、章の始まりなどでstyleが変わることがあるので微調整が必要ですが、基本はこれでOK

\hypersetup{
    colorlinks=true,       % 枠ではなく文字に色をつける
    linkcolor=blue!70!black, % 内部リンク（式番号など）の色
    citecolor=green!60!black, % 文献引用の色
    urlcolor=magenta,      % URLの色
    setpagesize=false      % 用紙サイズ設定の競合を防ぐ（jsarticle対策）
}


\renewcommand{\headfont}{\bfseries\rmfamily}
\renewcommand{\baselinestretch}{1.15}
\allowdisplaybreaks[1] % [1]～[4]で許容度を指定。1くらいが安全。
% 脚注を記号 (*, †, ‡, §, ¶...) にする
\renewcommand{\thefootnote}{\fnsymbol{footnote}}
\makesavenoteenv{derivationBox}

\numberwithin{equation}{section}

% derivationBoxの注釈関連
\usepackage{footmisc}
\renewcommand{\thefootnote}{\fnsymbol{footnote}}
\renewcommand{\thempfootnote}{\fnsymbol{footnote}}
% カウンタを共有
\makeatletter
\let\c@mpfootnote\c@footnote
\makeatother

% --- タイトル情報 ---
%タイトル
\title{ゼミノート \\ \large 物性物理のための場の理論・グリーン関数 \\ 小形正男} %タイトル
\author{田中響}%著者

\date{\today}  %任意の補足データ(\todayと入れれば今日の日付が出る)

%本文開始
\begin{document}
\maketitle %タイトル
\tableofcontents %目次

\maketitle

\newpage
\section{第二量子化}
\subsection{多体問題のSchrödinger方程式}
    同種N粒子系のSchrödinger方程式は
    \begin{align}
        &i\hbar \frac{\partial}{\partial t} \Psi(\bm{r}_1,\cdots,\bm{r}_N;t) = \hat{H} \Psi(\bm{r}_1,\cdots,\bm{r}_N;t) \label{NryuShiSyuredi}\\
        &\hat{H} = \sum_{i=1}^N \left( -\frac{\hbar^2}{2m_i} \nabla_i^2 + V_i(\bm{r}_i) \right) + \frac{1}{2}\sum_{i\ne j}U_{ij}(\bm{r}_i-\bm{r}_j)
    \end{align}
    となる。同種粒子波動関数$\Psi$は以下のいずれかの関係式を満たす
    \begin{align}
        \Psi(\bm{r}_1,\cdots, \bm{r}_i,\cdots, \bm{r}_j,\cdots,\bm{r}_N;t) &= \Psi(\bm{r}_1,\cdots, \bm{r}_j,\cdots, \bm{r}_i,\cdots,\bm{r}_N;t) \label{BoseTaisyouKansu} \\
        \Psi(\bm{r}_1,\cdots, \bm{r}_i,\cdots \bm{r}_j,\cdots,\bm{r}_N;t) &= -\Psi(\bm{r}_1,\cdots, \bm{r}_j,\cdots, \bm{r}_i,\cdots,\bm{r}_N;t)\label{FermiTaisyouKansu}
    \end{align}
    前者の式に従う粒子をBoson、後者をFermionという。
\subsection{完全対象関数と完全反対称関数}
    \eqref{BoseTaisyouKansu}、\eqref{FermiTaisyouKansu}式を満たす波動関数を表現するために、置換という考えを導入する。$\{1,2,\cdots,N\}$の数列を並び替えた数列を$\{P1,P2,\cdots,PN \}$と書くことにする。このような並び替えはすべてで$N!$通りある。これらすべての並び替えを足し合わせる記号を$\sum_P$とするとBosonの波動関数は
    \begin{align}
        \Psi_{\textrm{Bose}}(\bm{r}_1,\cdots,\bm{r}_N) = \sum_{P} \Psi(\bm{r}_{P1},\cdots,\bm{r}_{PN})
    \end{align}
    と書ける。\footnote{ここでは波動関数の規格化定数は無視する}Fermionの波動関数は
    \begin{align}
        \Psi_{\textrm{Fermi}}(\bm{r}_1,\cdots,\bm{r}_N) = \sum_{P} (-1)^P\Psi(\bm{r}_{P1},\cdots,\bm{r}_{PN})
    \end{align}
    となる。ここで$(-1)^P$のPは$\{1,2,\cdots,N\}$を$\{ P1,P2,\cdots,PN \}$に並び替える際に何回の数字の入れ替えを行ったかを表す値である。\\
    $\ \ $相互作用がない場合、一体のSchrödinger方程式
    \begin{align}
        \left\{ -\frac{\hbar^2}{2m} \nabla^2 + V(\bm{r})  \right\} \phi_n(\bm{r}) = \varepsilon_n \phi_n(\bm{r})
    \end{align}
    を満たす波動関数$\phi_n(\bm{r})$を用いると、相互作用がないN粒子ハミルトニアン
    \begin{align}
        \hat{H}_0 = \sum_{i=1}^N \left( -\frac{\hbar^2}{2m_i} \nabla_i^2 + V_i(\bm{r}_i) \right)
    \end{align}
    の固有状態は
    \begin{align}
        \phi_1(\bm{r}_1)\phi_2(\bm{r}_2)\cdots \phi_N(\bm{r}_N)
    \end{align}
    と書ける。したがってFermi粒子の波動関数は
    \begin{align}
        \Psi_{\textrm{Fermi}}(\bm{r}_1,\cdots,\bm{r}_N) = \sum_{P}(-1)^P \phi_1(\bm{r}_{P1})\phi_2(\bm{r}_{P2})\cdots\phi_N(\bm{r}_{PN})
    \end{align}
    となり、さらに行列式の定義より
    \begin{align}
        =\textrm{det} 
        \left\{
        \begin{matrix}
            \phi_1(\bm{r}_1)&\phi_2(\bm{r}_1)&\phi_3(\bm{r}_1)&\cdots&\phi_N(\bm{r}_1)\\
            \phi_1(\bm{r}_2)&\phi_2(\bm{r}_2)&\phi_3(\bm{r}_2)&\cdots&\phi_N(\bm{r}_2)\\
            \vdots&\vdots&\vdots&&\vdots\\
            \phi_1(\bm{r}_N)&\phi_2(\bm{r}_N)&\phi_3(\bm{r}_N)&\cdots&\phi_N(\bm{r}_N)
        \end{matrix}
        \right\}
    \end{align}
    とも表せる。この行列式をスレーター行列式という。
\subsection{第二量子化への準備}
    以下の交換関係を満たす演算子を定義する
    \begin{definition}[場の生成・消滅演算子]
        \begin{equation}
            \begin{aligned}
                \{ \hat{\psi}(\bm{r}),\hat{\psi}(\bm{r}') \} &= 0\\
                \{ \hat{\psi}^\dagger(\bm{r}),\hat{\psi}^\dagger(\bm{r}') \} &= 0\\
                \{ \hat{\psi}(\bm{r}),\hat{\psi}^\dagger(\bm{r}') \} &= \delta(\bm{r}-\bm{r}') \label{BanoKoukanKeiFermi}
            \end{aligned}
        \end{equation}
    \end{definition}
    粒子が全くいない真空を$\ket{0}$とすると、この演算子を用いて一つの粒子がいる状態を
    \begin{align}
        \hat{\psi}^\dagger(\bm{r}_1)\ket{0}
    \end{align}
    と書くことにする。また二つの粒子がいる状態を
    \begin{align}
        \frac{1}{\sqrt{2}}\hat{\psi}^\dagger(\bm{r}_1)\hat{\psi}^\dagger(\bm{r}_2)\ket{0}
    \end{align}
    と書け、N個の粒子がりる状態は
    \begin{align}
        \ket{\bm{r}_1,\cdots,\bm{r}_N}:=\frac{1}{\sqrt{N!}}\hat{\psi}^\dagger(\bm{r}_1)\cdots\hat{\psi}^\dagger(\bm{r}_N)\ket{0} \label{JyouTaiBectorK}
    \end{align}
    と書ける。この状態ベクトルのブラベクトルは
    \begin{align}
        \bra{\bm{r}_1,\cdots,\bm{r}_N} = \frac{1}{\sqrt{N!}}\bra{0}\hat{\psi}(\bm{r}_N)\cdots\hat{\psi}(\bm{r}_1)\label{JyouTaiBectorB}
    \end{align}
    となる。この状態ベクトルに対して以下が成り立つ
    \begin{formula}[状態ベクトルの内積]
        \begin{equation}
            \braket{\bm{r}_1,\cdots,\bm{r}_N|{\bm{r}'_1,\cdots,\bm{r}'_{N'}}} = \frac{\delta_{NN'}}{N!}\sum_P (-1)^P \delta(\bm{r}_1-\bm{r}'_{P1})\cdots\delta(\bm{r}_N-\bm{r}'_{PN}) \label{JyouTaiNaisek}
        \end{equation}
    \end{formula}
    \begin{derivationBox}
        \eqref{BanoKoukanKeiFermi}式の交換関係から得られる式
        \begin{align}
            \hat{\psi}(\bm{r})\hat{\psi}^\dagger(\bm{r})=-\hat{\psi}^\dagger(\bm{r})\hat{\psi}(\bm{r}) + \delta(\bm{r}-\bm{r}')
        \end{align}
        を用いると
        \begin{align}
            &\hat{\psi}(\bm{r}'_1) \hat{\psi}^\dagger(\bm{r}_1)\cdots\hat{\psi}(\bm{r}_N) \notag \\
            &= (-\hat{\psi}^\dagger(\bm{r}_1)\hat{\psi}(\bm{r}_1') + \delta(\bm{r}_1-\bm{r}'_1))\hat{\psi}^\dagger(\bm{r}_2)\cdots\hat{\psi}^\dagger(\bm{r}_N) \notag \\
            &= -\hat{\psi}^\dagger(\bm{r}_1)\hat{\psi}(\bm{r}_1') \hat{\psi}^\dagger(\bm{r}_2)\cdots\hat{\psi}^\dagger(\bm{r}_N) +  \delta(\bm{r}_1-\bm{r}'_1)\hat{\psi}^\dagger(\bm{r}_2)\cdots\hat{\psi}^\dagger(\bm{r}_N) \notag  \\
            &= -\hat{\psi}^\dagger(\bm{r}_1)( -\hat{\psi}^\dagger(\bm{r}_2)\hat{\psi}(\bm{r}_1') + \delta(\bm{r}_2-\bm{r}'_1))\hat{\psi}^\dagger(\bm{r}_3)\cdots\hat{\psi}^\dagger(\bm{r}_N) +  \delta(\bm{r}_1-\bm{r}'_1)\hat{\psi}^\dagger(\bm{r}_2)\cdots\hat{\psi}^\dagger(\bm{r}_N) \notag \\
            &= (-1)^2\hat{\psi}^\dagger(\bm{r}_1)\hat{\psi}^\dagger(\bm{r}_2)\hat{\psi}(\bm{r}_1') \hat{\psi}^\dagger(\bm{r}_3)\cdots\hat{\psi}^\dagger(\bm{r}_N) \notag \\
            &\ \ \ \ +(-1)\hat{\psi}^\dagger(\bm{r}_1)\delta(\bm{r}_2-\bm{r}'_1)\hat{\psi}^\dagger(\bm{r}_3)\cdots\hat{\psi}^\dagger(\bm{r}_N) + (-1)^0 \delta(\bm{r}_1-\bm{r}'_1)\hat{\psi}^\dagger(\bm{r}_2)\cdots\hat{\psi}^\dagger(\bm{r}_N) \notag \\
            &= (-1)^N \hat{\psi}^\dagger(\bm{r}_1)\cdots \hat{\psi}^\dagger(\bm{r}_N)\hat{\psi}(\bm{r}'_1) + \sum_{i=1}^{N} \left\{ \prod_{j=1}^{i-1} \hat{\psi}^\dagger(\bm{r}_j)\right\}\delta(\bm{r}_i - \bm{r}'_1) \left\{ \prod_{k=i+1}^{N} \hat{\psi}^\dagger(\bm{r}_k) \right\}
        \end{align}
        となる。第一項目は右端に消滅演算子があるため後に真空ケットに作用させると0となる。第二項目は数列$\{1,\cdots,N\}$を1に関しての並び替えを全て足し合わせ、元々$\hat{\psi}^\dagger(\bm{r}_1)$の代わりに$\delta(\bm{r}_1-\bm{r}'_1)$にしたものとなって、元の数列からの置換の回数に応じて$-1$の係数がかけられている。さらに左端に$\hat{\psi}(\bm{r}_2)$を作用させると、数列$\{1,\cdots,N\}$を1,2に関しての並び替えを全て足し合わせたものが得られる。したがって$\hat{psi}(\bm{r}_1),\cdots,\hat{\psi}(\bm{r}_N)$を順に左端から作用させると1~Nまでの並び替えを全て足したものが得られるので
        \begin{align}
            &\hat{\psi}(\bm{r}_N)\cdots\hat{\psi}(\bm{r}_1) \hat{\psi}^\dagger(\bm{r}'_1)\cdots\hat{\psi}^\dagger(\bm{r}'_{N'})\ket{0} \notag \\
            &= \delta_{NN'}\ket{0}\sum_P (-1)^P \delta(\bm{r}_1-\bm{r}'_{P1})\cdots \delta(\bm{r}_N-\bm{r}'_{P1})
        \end{align}
        となり\footnote{$N\ne N'$の場合は計算を行うと、生成演算子か消滅演算子のどちらかがあまり、それらが直接状態ベクトルに作用するので0となる}、\eqref{JyouTaiBectorK}、\eqref{JyouTaiBectorB}の定義式より
        \begin{align}
            \braket{\bm{r}_1,\cdots,\bm{r}_N|{\bm{r}'_1,\cdots,\bm{r}'_N}} = \frac{\delta_{NN'}}{N!}\sum_P (-1)^P \delta(\bm{r}_1-\bm{r}'_{P1})\cdots\delta(\bm{r}_N-\bm{r}'_{PN})
        \end{align}
        が成り立つ。
    \end{derivationBox}
\subsection{第二量子化の定式化}
    \begin{definition}[波動関数の状態ベクトル]
        \begin{equation}
            \ket{\Psi_N(t)} = \int d\bm{r}_1\cdots d\bm{r}_N\ \Psi(\bm{r}_1,\cdots,\bm{r}_N) \ket{\bm{r}_1, \cdots ,\bm{r}_N}
        \end{equation}
    \end{definition}
    この波動関数と、\eqref{JyouTaiBectorK}式との内積をとると\eqref{JyouTaiNaisek}より
    \begin{align}
        \braket{\bm{r}_1,\cdots ,\bm{r}_N| \Psi_N(t)} &= \bra{\bm{r}_1,\cdots ,\bm{r}_N} \int d\bm{r}'_1\cdots d\bm{r}'_N\ \Psi(\bm{r}'_1,\cdots,\bm{r}'_N) \ket{\bm{r}'_1, \cdots ,\bm{r}'_N} \notag \\
        &= \int d\bm{r}'_1\cdots d\bm{r}'_N\ \Psi(\bm{r}'_1,\cdots,\bm{r}'_N)  \braket{\bm{r}_1,\cdots ,\bm{r}_N|\bm{r}'_1, \cdots ,\bm{r}'_N} \notag \\
        &= \int d\bm{r}'_1\cdots d\bm{r}'_N\ \Psi(\bm{r}'_1,\cdots,\bm{r}'_N) \left\{ \frac{\delta_{NN'}}{N!}\sum_P (-1)^P \delta(\bm{r}_1-\bm{r}'_{P1})\cdots\delta(\bm{r}_N-\bm{r}'_{PN}) \right\} \notag \\
        &= \int d\bm{r}'_1\cdots d\bm{r}'_N\ \Psi(\bm{r}'_1,\cdots,\bm{r}'_N) \left\{ \frac{\delta_{NN'}}{N!}\sum_P (-1)^P \delta(\bm{r}_{P1}-\bm{r}'_1)\cdots\delta(\bm{r}_{PN}-\bm{r}'_N) \right\} \notag \\
        &=\frac{1}{N!} \sum_{P} (-1)^P \Psi(\bm{r}_{P1},\cdots,\bm{r}_{PN}) =: \Psi_{\textrm{AS}}(\bm{r}_1,\cdots,\bm{r}_N)
    \end{align}
    となり、Fermi粒子の反対称性をもった波動関数が得られる。\footnote{ASはantisymmertrizedの略記である}\\
    $\ \ $次に波動関数の状態ベクトルがSchrödinger方程式
    \begin{align}
        i\hbar \frac{\partial}{\partial t}\ket{\Psi_N(t)} = \hat{\mathscr{H}} \ket{\Psi_N(t)} \label{JyouTaiSyureding}
    \end{align}
    に従うとしたとき、ハミルトニアン$\mathscr{H}$はどのように置けばよいのかを考える。ハミルトニアン$\hat{\mathscr{H}}$を一粒子ハミルトニアン$\hat{\mathscr{H}}_0$と相互作用ハミルトニアン$\hat{\mathscr{H}}_{\text{int}}$の和で書くと以下で与えられる。
    \begin{definition}[第二量子化ハミルトニアン]
        \begin{equation}
            \begin{aligned}
                \hat{\mathscr{H}}_0 &= \int d\bm{r}\left\{ \frac{\hbar^2}{2m} \bm{\nabla}\hat{\psi}^\dagger(\bm{r})\cdot \bm{\nabla}\hat{\psi}(\bm{r}) + V(\bm{r}) \hat{\psi}^\dagger(\bm{r})\hat{\psi}(\bm{r}) \right\}\\
                \hat{\mathscr{H}}_{\text{int}} &= \frac{1}{2} \int \int d\bm{r}d\bm{r}'\ \hat{\psi}^\dagger (\bm{r})\hat{\psi}(\bm{r}') U(\bm{r}-\bm{r}')\hat{\psi}(\bm{r}')\hat{\psi}(\bm{r})
            \end{aligned}
            \label{DaiNIryouSikaHmiru}
        \end{equation}
    \end{definition}
    \begin{derivationBox}
        ハミルトニアンが\eqref{DaiNIryouSikaHmiru}式で与えられ、波動関数の状態ベクトルが\eqref{JyouTaiSyureding}式に従うと仮定すると、波動関数のSchrödinger方程式が再現されること確認する。\\
        $\ \ $まずは簡単な一粒子ポテンシャル項から計算をする
        \begin{align}
            &\bra{\bm{r}'_1,\cdots,\bm{r}'_N}\int d\bm{r}\ V(\bm{r}) \hat{\psi}^\dagger(\bm{r})\hat{\psi}(\bm{r}) \ket{\Psi_N(t)}\notag \\
            &= \bra{\bm{r}'_1,\cdots,\bm{r}'_N}\int d\bm{r}\ V(\bm{r}) \hat{\psi}^\dagger(\bm{r})\hat{\psi}(\bm{r}) \left\{ \int d\bm{r}_1\cdots d\bm{r}_N\ \Psi(\bm{r}_1,\cdots,\bm{r}_N) \ket{\bm{r}_1, \cdots ,\bm{r}_N} \right\} \notag \\
            &= \int d\bm{r}\ \int d\bm{r}_1\cdots d\bm{r}_N\  V(\bm{r}) \Psi(\bm{r}_1,\cdots,\bm{r}_N) \bra{\bm{r}'_1,\cdots,\bm{r}'_N}\hat{\psi}^\dagger(\bm{r})\hat{\psi}(\bm{r})\ket{\bm{r}_1, \cdots ,\bm{r}_N} 
        \end{align}
        となり、ここで状態ベクトルの定義式\eqref{JyouTaiBectorK}より
        \begin{align}
            \hat{\psi}^\dagger(\bm{r})\hat{\psi}(\bm{r})\ket{\bm{r}_1, \cdots ,\bm{r}_N} 
            &= \hat{\psi}^\dagger(\bm{r})\hat{\psi}(\bm{r})\left\{ \frac{1}{\sqrt{N!}}\hat{\psi}^\dagger(\bm{r}_1)\cdots\hat{\psi}^\dagger(\bm{r}_N)\ket{0} \right\} \notag \\
            &= \frac{1}{\sqrt{N!}}\hat{\psi}^\dagger(\bm{r})\hat{\psi}(\bm{r}) \hat{\psi}^\dagger(\bm{r}_1)\cdots\hat{\psi}^\dagger(\bm{r}_N)\ket{0}
        \end{align}
        となり、ここで
        \begin{align}
            [ \hat{\psi}^\dagger(\bm{r}) \hat{\psi}(\bm{r}), \hat{\psi}^\dagger(\bm{r}')] &=  \hat{\psi}^\dagger(\bm{r}) \hat{\psi}(\bm{r})\hat{\psi}^\dagger(\bm{r}') - \hat{\psi}^\dagger(\bm{r}')\hat{\psi}^\dagger(\bm{r}) \hat{\psi}(\bm{r}) \notag \\
            &=  \hat{\psi}^\dagger(\bm{r}) \left\{ -\hat{\psi}^\dagger(\bm{r}')\hat{\psi}(\bm{r}) + \delta(\bm{r}-\bm{r'}) \right\} - \hat{\psi}^\dagger(\bm{r}')\hat{\psi}^\dagger(\bm{r}) \hat{\psi}(\bm{r}) \notag \\
            &=  \hat{\psi}^\dagger(\bm{r}')\hat{\psi}^\dagger(\bm{r})\hat{\psi}(\bm{r}) + \hat{\psi}^\dagger(\bm{r}) \delta(\bm{r}-\bm{r'}) - \hat{\psi}^\dagger(\bm{r}')\hat{\psi}^\dagger(\bm{r}) \hat{\psi}(\bm{r})\notag \\
            &=\hat{\psi}^\dagger(\bm{r}) \delta(\bm{r}-\bm{r'}) =\hat{\psi}^\dagger(\bm{r}') \delta(\bm{r}-\bm{r'}) 
        \end{align}
        の結果を用いると
        \begin{align}
            \frac{1}{\sqrt{N!}}\hat{\psi}^\dagger(\bm{r})\hat{\psi}(\bm{r}) \hat{\psi}^\dagger(\bm{r}_1)\cdots\hat{\psi}^\dagger(\bm{r}_N)\ket{0} &= \frac{1}{\sqrt{N!}}\left\{ \hat{\psi}^\dagger(\bm{r}_1)\hat{\psi}^\dagger(\bm{r})\hat{\psi}(\bm{r}) + \hat{\psi}^\dagger(\bm{r}) \delta(\bm{r}-\bm{r}_1)  \right\}\cdots\hat{\psi}^\dagger(\bm{r}_N)\ket{0} \notag \\
            &= \frac{1}{\sqrt{N!}} \hat{\psi}^\dagger(\bm{r}_1)\ \hat{\psi}^\dagger(\bm{r})\hat{\psi}(\bm{r})\ \hat{\psi}^\dagger(\bm{r}_2)\cdots\hat{\psi}^\dagger(\bm{r}_N)\ket{0} \notag \\
            &\ \ \ \ + \frac{1}{\sqrt{N!}} \delta(\bm{r}-\bm{r}_1)\hat{\psi}^\dagger(\bm{r}_1)\hat{\psi}^\dagger(\bm{r}_2)\cdots\hat{\psi}^\dagger(\bm{r}_N)\ket{0} \notag \\
            &= \frac{1}{\sqrt{N!}} \hat{\psi}^\dagger(\bm{r}_1)\ \hat{\psi}^\dagger(\bm{r})\hat{\psi}(\bm{r})\ \hat{\psi}^\dagger(\bm{r}_2)\cdots\hat{\psi}^\dagger(\bm{r}_N)\ket{0} \notag \\
            &\ \ \ \ + \delta(\bm{r}-\bm{r}_1)\ket{\bm{r}_1,\cdots,\bm{r}_N} \notag \\
            &= \frac{1}{\sqrt{N!}} \hat{\psi}^\dagger(\bm{r}_1)\hat{\psi}^\dagger(\bm{r}_2)\ \hat{\psi}^\dagger(\bm{r})\hat{\psi}(\bm{r})\ \hat{\psi}^\dagger(\bm{r}_3)\cdots\hat{\psi}^\dagger(\bm{r}_N)\ket{0} \notag \\
            &\ \ \ \ + \sum_{i=1}^2\delta(\bm{r}-\bm{r}_i)\ket{\bm{r}_1,\cdots,\bm{r}_N} \notag \\
            &= \frac{1}{\sqrt{N!}} \hat{\psi}^\dagger(\bm{r}_1)\cdots \hat{\psi}(\bm{r}_N) \hat{\psi}^\dagger(\bm{r})\hat{\psi}(\bm{r})\ket{0} \notag \\
            &\ \ \ \ + \sum_{i=1}^N\delta(\bm{r}-\bm{r}_i)\ket{\bm{r}_1,\cdots,\bm{r}_N} \notag \\
            &= \sum_{i=1}^N\delta(\bm{r}-\bm{r}_i)\ket{\bm{r}_1,\cdots,\bm{r}_N} \label{PotePoteKou}
        \end{align}
        となる。したがって
        \begin{align}
            &\int d\bm{r}\ \int d\bm{r}_1\cdots d\bm{r}_N\   V(\bm{r})\Psi(\bm{r}_1,\cdots,\bm{r}_N)\bra{\bm{r}'_1,\cdots,\bm{r}'_N}\hat{\psi}^\dagger(\bm{r})\hat{\psi}(\bm{r})\ket{\bm{r}_1, \cdots ,\bm{r}_N} \notag \\
            &= \int d\bm{r}\ \int d\bm{r}_1\cdots d\bm{r}_N\   V(\bm{r})\Psi(\bm{r}_1,\cdots,\bm{r}_N)\bra{\bm{r}'_1,\cdots,\bm{r}'_N} \left\{  \sum_{i=1}^N\delta(\bm{r}-\bm{r}_i)\ket{\bm{r}_1,\cdots,\bm{r}_N} \right\} \notag \\
            &= \sum_{i=1}^N \int d\bm{r}\ \int d\bm{r}_1\cdots d\bm{r}_N\   V(\bm{r}) \Psi(\bm{r}_1,\cdots,\bm{r}_N) \delta(\bm{r}-\bm{r}_i) \braket{\bm{r}'_1,\cdots,\bm{r}'_N| \bm{r}_1,\cdots,\bm{r}_N} \notag \\
            &= \sum_{i=1}^N \int d\bm{r}_1\cdots d\bm{r}_N\  V(\bm{r}_i) \Psi(\bm{r}_1,\cdots,\bm{r}_N)  \braket{\bm{r}'_1,\cdots,\bm{r}'_N| \bm{r}_1,\cdots,\bm{r}_N} \notag \\
            &= \sum_{i=1}^N \int d\bm{r}_1\cdots d\bm{r}_N\   V(\bm{r}_i)\Psi(\bm{r}_1,\cdots,\bm{r}_N)  \left\{ \frac{1}{N!}\sum_P (-1)^P \delta(\bm{r}_1-\bm{r}'_{P1})\cdots\delta(\bm{r}_N-\bm{r}'_{PN})  \right\} \notag \\
            &= \sum_{i=1}^N \left\{ \frac{1}{N!} \sum_{P}V(\bm{r}'_{Pi})(-1)^P \Psi(\bm{r}'_{P1},\cdots,\bm{r}'_{PN}) \right\} \notag \\
            &= \frac{1}{N!} \sum_{P}\left\{ \sum_{i=1}^NV(\bm{r}'_{Pi})\right\} (-1)^P \Psi(\bm{r}'_{P1},\cdots,\bm{r}'_{PN})  \notag \\
            &= \frac{1}{N!} \sum_{P}\left\{ \sum_{i=1}^NV(\bm{r}'_i)\right\} (-1)^P \Psi(\bm{r}'_{P1},\cdots,\bm{r}'_{PN})  \notag \\
            &= \left\{ \sum_{i=1}^NV(\bm{r}'_i)\right\} \frac{1}{N!} \sum_{P} (-1)^P \Psi(\bm{r}'_{P1},\cdots,\bm{r}'_{PN})  \notag \\
            &= \sum_{i=1}^N V(\bm{r}'_i) \Psi_{\text{AS}}(\bm{r}'_1,\cdots,\bm{r}'_N)
        \end{align}
        となる。\\
        $\ \ $次に運動エネルギー項を計算する。
        \begin{align}
            &\bra{\bm{r}'_1,\cdots,\bm{r}'_N}\int d\bm{r}\              \frac{\hbar^2}{2m}\bm{\nabla}\hat{\psi}^\dagger(\bm{r})\cdot \bm{\nabla}\hat{\psi}(\bm{r}) \ket{\Psi_N(t)}\notag \\
            &= \bra{\bm{r}'_1,\cdots,\bm{r}'_N}\int d\bm{r}\  \frac{\hbar^2}{2m}\bm{\nabla}\hat{\psi}^\dagger(\bm{r})\cdot \bm{\nabla}\hat{\psi}(\bm{r}) \left\{ \int d\bm{r}_1\cdots d\bm{r}_N\ \Psi(\bm{r}_1,\cdots,\bm{r}_N) \ket{\bm{r}_1, \cdots ,\bm{r}_N} \right\} \notag \\
            &= \frac{\hbar^2}{2m}\int d\bm{r} \int d\bm{r}_1\cdots d\bm{r}_N\   \Psi(\bm{r}_1,\cdots,\bm{r}_N) \bra{\bm{r}'_1,\cdots,\bm{r}'_N}\bm{\nabla}\hat{\psi}^\dagger(\bm{r})\cdot\bm{\nabla}\hat{\psi}(\bm{r})\ket{\bm{r}_1, \cdots ,\bm{r}_N} 
        \end{align}
        となる。ここで
        \begin{align}
            [\bm{\nabla}\hat{\psi}^\dagger(\bm{r})\cdot\bm{\nabla}\hat{\psi}(\bm{r}), \hat{\psi}^\dagger(\bm{r}')]&= \bm{\nabla}\hat{\psi}^\dagger(\bm{r})\cdot\bm{\nabla}\hat{\psi}(\bm{r}) \hat{\psi}^\dagger(\bm{r}') - \hat{\psi}^\dagger(\bm{r}')\bm{\nabla}\hat{\psi}^\dagger(\bm{r})\cdot\bm{\nabla}\hat{\psi}(\bm{r}) \notag \\
            &= \bm{\nabla}\hat{\psi}^\dagger(\bm{r})\cdot\bm{\nabla}\left\{- \hat{\psi}^\dagger(\bm{r}')\hat{\psi}(\bm{r}) +\delta(\bm{r}-\bm{r}') \right\} - \bm{\nabla}\left\{ -\hat{\psi}^\dagger(\bm{r})\hat{\psi}^\dagger(\bm{r}') \right\}\cdot\bm{\nabla}\hat{\psi}(\bm{r}) \notag \\
            &= -\bm{\nabla}\hat{\psi}^\dagger(\bm{r})\cdot\bm{\nabla} \hat{\psi}^\dagger(\bm{r}')\hat{\psi}(\bm{r}) +\bm{\nabla}\hat{\psi}^\dagger(\bm{r})\cdot\bm{\nabla}\delta(\bm{r}-\bm{r}')  +\bm{\nabla}\hat{\psi}^\dagger(\bm{r})\cdot\bm{\nabla}\hat{\psi}^\dagger(\bm{r}')\hat{\psi}(\bm{r}) \notag \\
            &= \bm{\nabla}\hat{\psi}^\dagger(\bm{r})\cdot\bm{\nabla}\delta(\bm{r}-\bm{r}') 
        \end{align}
        の結果を用いると、一粒子ポテンシャルのときと比べ、生成演算子とデルタ関数にナブラ演算子がつくことしか変わらないので、同様の計算により
        \begin{align}
            &\frac{1}{\sqrt{N!}}\bm{\nabla}\hat{\psi}^\dagger(\bm{r})\cdot \bm{\nabla}\hat{\psi}(\bm{r}) \hat{\psi}^\dagger(\bm{r}_1)\cdots\hat{\psi}^\dagger(\bm{r}_N)\ket{0} \notag \\
            &= \frac{1}{\sqrt{N!}}\left\{ \hat{\psi}^\dagger(\bm{r}_1)\bm{\nabla}\hat{\psi}^\dagger(\bm{r})\cdot\bm{\nabla}\hat{\psi}(\bm{r}) + \bm{\nabla}\hat{\psi}^\dagger(\bm{r})\cdot\bm{\nabla} \delta(\bm{r}-\bm{r}_1)  \right\}\cdots\hat{\psi}^\dagger(\bm{r}_N)\ket{0} \notag \\
            &= \frac{1}{\sqrt{N!}} \hat{\psi}^\dagger(\bm{r}_1)\ \bm{\nabla}\hat{\psi}^\dagger(\bm{r})\cdot \bm{\nabla}\hat{\psi}(\bm{r})\ \hat{\psi}^\dagger(\bm{r}_2)\cdots\hat{\psi}^\dagger(\bm{r}_N)\ket{0} \notag \\
            &\ \ \ \ + \frac{1}{\sqrt{N!}} \bm{\nabla}\delta(\bm{r}-\bm{r}_1)\cdot \bm{\nabla}\hat{\psi}^\dagger(\bm{r})\hat{\psi}^\dagger(\bm{r}_2)\cdots\hat{\psi}^\dagger(\bm{r}_N)\ket{0} \notag \\
            &= \frac{1}{\sqrt{N!}} \hat{\psi}^\dagger(\bm{r}_1)\ \bm{\nabla}\hat{\psi}^\dagger(\bm{r})\cdot \bm{\nabla}\hat{\psi}(\bm{r})\ \hat{\psi}^\dagger(\bm{r}_2)\cdots\hat{\psi}^\dagger(\bm{r}_N)\ket{0} \notag \\
            &\ \ \ \ + \bm{\nabla}\delta(\bm{r}-\bm{r}_1)\cdot \bm{\nabla}\ket{\bm{r},\bm{r}_2,\cdots,\bm{r}_N} \notag \\
            &= \frac{1}{\sqrt{N!}} \hat{\psi}^\dagger(\bm{r}_1)\cdots\hat{\psi}^\dagger(\bm{r}_N)\ \bm{\nabla}\hat{\psi}^\dagger(\bm{r})\cdot \bm{\nabla}\hat{\psi}(\bm{r})\ket{0} \notag \\
            &\ \ \ \ + \sum_{i=1}^N \bm{\nabla}\delta(\bm{r}-\bm{r}_i)\cdot \bm{\nabla}\ket{\bm{r},\cdots,\bm{r}_i{=\bm{r}},\cdots,\bm{r}_N} \notag \\
            &= \sum_{i=1}^N \bm{\nabla}\delta(\bm{r}-\bm{r}_i)\cdot \bm{\nabla}\ket{\bm{r},\cdots,\bm{r}_i{=\bm{r}},\cdots,\bm{r}_N}
        \end{align}
        となるので
        \begin{align}
            &\frac{\hbar^2}{2m}\int d\bm{r} \int d\bm{r}_1\cdots d\bm{r}_N\   \Psi(\bm{r}_1,\cdots,\bm{r}_N) \bra{\bm{r}'_1,\cdots,\bm{r}'_N}\bm{\nabla}\hat{\psi}^\dagger(\bm{r})\cdot\bm{\nabla}\hat{\psi}(\bm{r})\ket{\bm{r}_1, \cdots ,\bm{r}_N} \notag \\
            &=\frac{\hbar^2}{2m}\int d\bm{r} \int d\bm{r}_1\cdots d\bm{r}_N\   \Psi(\bm{r}_1,\cdots,\bm{r}_N) \bra{\bm{r}'_1,\cdots,\bm{r}'_N}\left\{ \sum_{i=1}^N \bm{\nabla}\delta(\bm{r}-\bm{r}_i)\cdot \bm{\nabla}\ket{\bm{r},\cdots,\bm{r}_i{=\bm{r}},\cdots,\bm{r}_N} \right\} \notag \\
            &= -\frac{\hbar^2}{2m}\int d\bm{r}\int d\bm{r}_1\cdots d\bm{r}_N\   \Psi(\bm{r}_1,\cdots,\bm{r}_N) \bra{\bm{r}'_1,\cdots,\bm{r}'_N}\left\{ \sum_{i=1}^N \delta(\bm{r}-\bm{r}_i) \bm{\nabla}^2\ket{\bm{r},\cdots,\bm{r}_i{=\bm{r}},\cdots,\bm{r}_N} \right\} \notag \\
            &= -\frac{\hbar^2}{2m}\int d\bm{r}_1\cdots d\bm{r}_N\   \Psi(\bm{r}_1,\cdots,\bm{r}_N) \bra{\bm{r}'_1,\cdots,\bm{r}'_N}\left\{ \sum_{i=1}^N \bm{\nabla}^2_{\bm{r}_i}\ket{\bm{r}_1,\cdots,\bm{r}_N} \right\} \notag \\
            &= \sum_{i=1}^N\int d\bm{r}_1\cdots d\bm{r}_N\   \left( -\frac{\hbar^2}{2m} \bm{\nabla}^2_{\bm{r}_i}\right)\Psi(\bm{r}_1,\cdots,\bm{r}_N) \braket{\bm{r}'_1,\cdots,\bm{r}'_N|\bm{r}_1,\cdots,\bm{r}_N} \notag \\
            &= \sum_{i=1}^N\int d\bm{r}_1\cdots d\bm{r}_N\   \left( -\frac{\hbar^2}{2m} \bm{\nabla}^2_{\bm{r}_i}\right)\Psi(\bm{r}_1,\cdots,\bm{r}_N) \left\{ \frac{1}{N!}\sum_P (-1)^P \delta(\bm{r}_1-\bm{r}'_{P1})\cdots\delta(\bm{r}_N-\bm{r}'_{PN})  \right\} \notag \\
            &= \sum_{i=1}^N \frac{1}{N!}\sum_P \left( -\frac{\hbar^2}{2m} \bm{\nabla}^2_{\bm{r}'_{Pi}}\right)(-1)^P \Psi(\bm{r}'_{P1},\cdots,\bm{r}'_{PN})  \notag \\
            &=  \frac{1}{N!}\sum_P \left\{ \sum_{i=1}^N\left( -\frac{\hbar^2}{2m} \bm{\nabla}^2_{\bm{r}'_{Pi}}\right)\right\} (-1)^P \Psi(\bm{r}'_{P1},\cdots,\bm{r}'_{PN})  \notag \\
            &=  \frac{1}{N!}\sum_P \left\{ \sum_{i=1}^N\left( -\frac{\hbar^2}{2m} \bm{\nabla}^2_{\bm{r}'_i}\right)\right\} (-1)^P \Psi(\bm{r}'_{P1},\cdots,\bm{r}'_{PN})  \notag \\
            &= \left\{ \sum_{i=1}^N\left( -\frac{\hbar^2}{2m} \bm{\nabla}^2_{\bm{r}'_i}\right)\right\} \Psi_{\text{AS}}(\bm{r}'_1,\cdots,\bm{r}_N)
        \end{align}
        となる。\\
        $\ \ $最後に相互作用項を計算する。
        \begin{align}
            &\bra{\bm{r}'_1,\cdots,\bm{r}'_N}\frac{1}{2}\int d\bm{r}\ U(\bm{r}-\bm{r}') \hat{\psi}^\dagger(\bm{r})\hat{\psi}^\dagger(\bm{r}')\hat{\psi}(\bm{r}')\hat{\psi}(\bm{r}) \ket{\Psi_N(t)}\notag \\
            &= \bra{\bm{r}'_1,\cdots,\bm{r}'_N}\frac{1}{2}\int d\bm{r}\ U(\bm{r}-\bm{r}') \hat{\psi}^\dagger(\bm{r})\hat{\psi}^\dagger(\bm{r}')\hat{\psi}(\bm{r}')\hat{\psi}(\bm{r})  \left\{ \int d\bm{r}_1\cdots d\bm{r}_N\ \Psi(\bm{r}_1,\cdots,\bm{r}_N) \ket{\bm{r}_1, \cdots ,\bm{r}_N} \right\} \notag \\
            &= \frac{1}{2}\int d\bm{r}\ \int d\bm{r}_1\cdots d\bm{r}_N\  U(\bm{r}-\bm{r}') \Psi(\bm{r}_1,\cdots,\bm{r}_N) \bra{\bm{r}'_1,\cdots,\bm{r}'_N}\hat{\psi}^\dagger(\bm{r})\hat{\psi}^\dagger(\bm{r}')\hat{\psi}(\bm{r}')\hat{\psi}(\bm{r}) \ket{\bm{r}_1, \cdots ,\bm{r}_N} 
        \end{align}
        となり
        \begin{align}
            &\hat{\psi}^\dagger(\bm{r})\hat{\psi}^\dagger(\bm{r}')\hat{\psi}(\bm{r}')\hat{\psi}(\bm{r}) \ket{\bm{r}_1, \cdots ,\bm{r}_N} \notag \\
            &= \hat{\psi}^\dagger(\bm{r})\hat{\psi}^\dagger(\bm{r}')\hat{\psi}(\bm{r}')\sum_{i=1}^N(-1)^{i-1}\delta(\bm{r}-\bm{r_i}) \ket{\bm{r}_1,\cdots,{\color{lightgray}{\bm{r}_i}}, \cdots ,\bm{r}_N} \notag \\
            &= \hat{\psi}^\dagger(\bm{r})\hat{\psi}^\dagger(\bm{r}')2\sum_{j> i}(-1)^{i-1+j-2}\delta(\bm{r}-\bm{r_i}) \delta(\bm{r}'-\bm{r_j}) \ket{\bm{r}_1,\cdots,{\color{lightgray}{\bm{r}_i}},\cdots,{\color{lightgray}{\bm{r}_j}}, \cdots ,\bm{r}_N}  \notag \\
            &= \hat{\psi}^\dagger(\bm{r})2\sum_{j> i}(-1)^{i-1+j-2+j-2}\delta(\bm{r}-\bm{r_i}) \delta(\bm{r}'-\bm{r_j}) \ket{\bm{r}_1,\cdots,{\color{lightgray}{\bm{r}_i}},\cdots,{\color{lightgray}{\bm{r}_j}}, \cdots ,\bm{r}_N} \notag \\
            &= 2\sum_{j> i}(-1)^{i-1+j-2+j-2+ i-1}\delta(\bm{r}-\bm{r_i}) \delta(\bm{r}'-\bm{r_j}) \ket{\bm{r}_1,\cdots,{\bm{r}},\cdots,\bm{r}_j, \cdots ,\bm{r}_N} \notag \\
            &= 2\sum_{j> i}\delta(\bm{r}-\bm{r_i}) \delta(\bm{r}'-\bm{r_j}) \ket{\bm{r}_1,\cdots,\bm{r}_i=\bm{r},\cdots,\bm{r}_j=\bm{r}', \cdots ,\bm{r}_N} 
        \end{align}
        となるので
        \begin{align}
            &\frac{1}{2}\int d\bm{r}\ \int d\bm{r}_1\cdots d\bm{r}_N\  U(\bm{r}-\bm{r}') \Psi(\bm{r}_1,\cdots,\bm{r}_N) \bra{\bm{r}'_1,\cdots,\bm{r}'_N}\hat{\psi}^\dagger(\bm{r})\hat{\psi}^\dagger(\bm{r}')\hat{\psi}(\bm{r}')\hat{\psi}(\bm{r}) \ket{\bm{r}_1, \cdots ,\bm{r}_N}  \notag \\
            &=\frac{1}{2}\int d\bm{r}\ \int d\bm{r}_1\cdots d\bm{r}_N\  U(\bm{r}-\bm{r}') \Psi(\bm{r}_1,\cdots,\bm{r}_N) \notag \\
            &\ \ \ \ \ \ \ \ \bra{\bm{r}'_1,\cdots,\bm{r}'_N}\left\{ 2\sum_{j> i}\delta(\bm{r}-\bm{r_i}) \delta(\bm{r}'-\bm{r_j}) \ket{\bm{r}_1,\cdots,\bm{r}_i=\bm{r},\cdots,\bm{r}_j=\bm{r}', \cdots ,\bm{r}_N}  \right\} \notag \\
            &=\sum_{j> i}\int d\bm{r}\ \int d\bm{r}_1\cdots d\bm{r}_N\  U(\bm{r}-\bm{r}') \Psi(\bm{r}_1,\cdots,\bm{r}_N) \notag \\
            &\ \ \ \ \ \ \ \  \delta(\bm{r}-\bm{r_i}) \delta(\bm{r}'-\bm{r_j}) \braket{\bm{r}'_1,\cdots,\bm{r}'_N|\bm{r}_1,\cdots,\bm{r}_i=\bm{r},\cdots,\bm{r}_j=\bm{r}', \cdots ,\bm{r}_N}  \notag \\
            &=\sum_{j> i} \int d\bm{r}_1\cdots d\bm{r}_N\  U(\bm{r}_i-\bm{r}_j) \Psi(\bm{r}_1,\cdots,\bm{r}_N)\braket{\bm{r}'_1,\cdots ,\bm{r}'_N|\bm{r}_1,\cdots,\bm{r}_N}  \notag \\
            &=\sum_{j> i} \int d\bm{r}_1\cdots d\bm{r}_N\  U(\bm{r}_i-\bm{r}_j) \Psi(\bm{r}_1,\cdots,\bm{r}_N)\left\{ \frac{1}{N!}\sum_P (-1)^P \delta(\bm{r}_1-\bm{r}'_{P1})\cdots\delta(\bm{r}_N-\bm{r}'_{PN}) \right\} \notag \\
            &=\sum_{j> i} \frac{1}{N!}\sum_P U(\bm{r}_{Pi}-\bm{r}_{Pj}) (-1)^P \Psi(\bm{r}_{P1},\cdots,\bm{r}_{PN})  \notag \\
            &= \frac{1}{N!}\sum_P \left\{ \sum_{j> i}U(\bm{r}_{Pi}-\bm{r}_{Pj}) \right\} (-1)^P \Psi(\bm{r}_{P1},\cdots,\bm{r}_{PN})  \notag \\
            &= \frac{1}{N!}\sum_P \left\{ \sum_{j> i}U(\bm{r}_{i}-\bm{r}_{j}) \right\} (-1)^P \Psi(\bm{r}_{P1},\cdots,\bm{r}_{PN}) \notag \\
            &= \left\{ \sum_{j> i}U(\bm{r}_{i}-\bm{r}_{j}) \right\}\Psi_{\text{AS}}(\bm{r}_1,\cdots,\bm{r}_N)
        \end{align}
        となる。\\
        $\ \ $したがって
        \begin{align}
            \bra{\bm{r}'_1,\cdots,\bm{r}'_N} \hat{\mathscr{H}} \ket{\Psi_N(t)} = \left[ \sum_{i=1}^N \left\{ -\frac{\hbar^2}{2m}\nabla^2_{\bm{r}_i} + V(\bm{r}_i) \right\} + \sum_{j> i}U(\bm{r}_{i}-\bm{r}_{j}) \right]\Psi_{\text{AS}}(\bm{r}_1,\cdots,\bm{r}_N) 
        \end{align}
        となり、Schrödinger方程式が再現される。
    \end{derivationBox}
\subsection{物理量の期待値}
    系のエネルギー$\hat{H}$の期待値は、波動関数を用いて
    \begin{align}
        \braket{\hat{H}} &= \int d\bm{r}_1 \cdots d\bm{r}_N\ \Psi_{\text{AS}}^* (\bm{r}_1,\cdots,\bm{r}_N) \hat{H}\Psi_{\text{AS}} (\bm{r}_1,\cdots,\bm{r}_N) \\
        &\hat{H} = \sum_{i=1}^N \left\{ -\frac{\hbar^2}{2m}\nabla^2_{\bm{r}_i} + V(\bm{r}_i) \right\} + \sum_{j> i}U(\bm{r}_{i}-\bm{r}_{j})
    \end{align}
    で表される。前節で導入した波動関数の状態ベクトルを用いて、この期待値がどのように表されるのかを考える
    \begin{formula}[エネルギーの期待値]
        \begin{equation}
            \braket{\hat{H}} = \braket{\Psi_N(t)|\hat{\mathscr{H}}|\Psi_N(t)}
        \end{equation}
    \end{formula}
    \begin{derivationBox}
        \begin{align}
            \braket{\hat{H}} &= \int d\bm{r}_1 \cdots d\bm{r}_N\ \Psi_{\text{AS}}^* (\bm{r}_1,\cdots,\bm{r}_N)\hat{H}\Psi_{\text{AS}} (\bm{r}_1,\cdots,\bm{r}_N) \notag \\
            &= \int d\bm{r}_1 \cdots d\bm{r}_N\ \Psi_{\text{AS}}^* (\bm{r}_1,\cdots,\bm{r}_N)\bra{\bm{r}_1,\cdots,\bm{r}_N} \hat{\mathscr{H}} \ket{\Psi_N(t)} \notag \\ 
            &= \int d\bm{r}_1 \cdots d\bm{r}_N\ \braket{\Psi_N(t)|\bm{r}_1,\cdots,\bm{r}_N}\bra{\bm{r}_1,\cdots,\bm{r}_N} \hat{\mathscr{H}} \ket{\Psi_N(t)} \notag \\
            &= \braket{\Psi_N(t)|\hat{\mathscr{H}}|\Psi_N(t)}
        \end{align}
    \end{derivationBox}
    位置$r$での粒子数密度の期待値は
    \begin{align}
        \braket{\hat{\bm{r}}} &=\int d\bm{r}_1 \cdots d\bm{r}_N\ \Psi_{\text{AS}}^* (\bm{r}_1,\cdots,\bm{r}_N) \hat{\bm{r}}\Psi_{\text{AS}} (\bm{r}_1,\cdots,\bm{r}_N) \\
        &\hat{\bm{r}} = \sum_{i=1}^N \delta(\bm{r}-\bm{r}_i)
    \end{align}
    と表される。
    \begin{formula}[粒子数密度の期待値]
        \begin{equation}
            \braket{\hat{\bm{r}}} = \braket{\Psi_N(t)|\hat{\rho}(\bm{r})|\Psi_N(t)}
        \end{equation}
        ここで密度演算子$\hat{\rho}(\bm{r})$は
        \begin{equation}
            \hat{\rho}(\bm{r}):=\hat{\psi}^\dagger(\bm{r})\hat{\psi}(\bm{r})
        \end{equation}
        と定義した
    \end{formula}
    \begin{derivationBox}
        \eqref{PotePoteKou}式の計算を思い出すと
        \begin{align}
            &\hat{\psi}^\dagger(\bm{r})\hat{\psi}(\bm{r}) \ket{\bm{r}_1,\cdots,\bm{r}_N} \notag \\
            &=
            \hat{\psi}^\dagger(\bm{r})\hat{\psi}(\bm{r})\ \hat{\psi}^\dagger(\bm{r}_1)\cdots\hat{\psi}^\dagger(\bm{r}_N)\ket{0} \notag \\
            &= \sum_{i=1}^N\delta(\bm{r}-\bm{r}_i)\ket{\bm{r}_1,\cdots,\bm{r}_N}
        \end{align}
        となるので
        \begin{align}
            \bra{\bm{r}_1,\cdots,\bm{r}_N}\hat{\psi}^\dagger(\bm{r})\hat{\psi}(\bm{r}) \ket{\Psi_N(t)} &= \bra{\bm{r}_1,\cdots,\bm{r}_N}\int d\bm{r}_1\cdots d\bm{r}_N\ \Psi_N(\bm{r}_1,\cdots,\bm{r}_N) \hat{\psi}^\dagger(\bm{r})\hat{\psi}(\bm{r})\ket{\bm{r}_1,\cdots,\bm{r}_N} \notag  \\
            &= \bra{\bm{r}_1,\cdots,\bm{r}_N}\int d\bm{r}_1\cdots d\bm{r}_N\ \Psi_N(\bm{r}_1,\cdots,\bm{r}_N) \sum_{i=1}^N\delta(\bm{r}-\bm{r}_i)\ket{\bm{r}_1,\cdots,\bm{r}_N} \notag \\
            &= \sum_{i=1}^N\delta(\bm{r}-\bm{r}_i) \braket{\bm{r}_1,\cdots,\bm{r}_N|\Psi_N(t)} \notag \\
            &= \sum_{i=1}^N\delta(\bm{r}-\bm{r}_i) \Psi_{\text{AS}}(\bm{r}_1,\cdots,\bm{r}_N)
        \end{align}
        が成り立ち、これを用いると
        \begin{align}
            \braket{\hat{\bm{r}}} &=\int d\bm{r}_1 \cdots d\bm{r}_N\ \Psi_{\text{AS}}^* (\bm{r}_1,\cdots,\bm{r}_N) \sum_{i=1}^N \delta(\bm{r}-\bm{r}_i)\Psi_{\text{AS}} (\bm{r}_1,\cdots,\bm{r}_N) \\
        &= \int d\bm{r}_1 \cdots d\bm{r}_N\ \Psi_{\text{AS}}^* (\bm{r}_1,\cdots,\bm{r}_N) \bra{\bm{r}_1,\cdots,\bm{r}_N}\hat{\psi}^\dagger(\bm{r})\hat{\psi}(\bm{r}) \ket{\Psi_N(t)} \\
        &= \int d\bm{r}_1 \cdots d\bm{r}_N\ \braket{\Psi_N(t)|\bm{r}_1,\cdots,\bm{r}_N} \bra{\bm{r}_1,\cdots,\bm{r}_N}\hat{\psi}^\dagger(\bm{r})\hat{\psi}(\bm{r}) \ket{\Psi_N(t)} \\
        &=\bra{\Psi_N(t)}\hat{\psi}^\dagger(\bm{r})\hat{\psi}(\bm{r}) \ket{\Psi_N(t)} 
        \end{align}
        となる。
    \end{derivationBox}
\subsection{固有関数による展開}
    相互作用がない場合、一体のSchrödinger方程式の固有関数$\phi_n(\bm{r})$が求まったとする。この固有関数で場の演算子を展開する
    \begin{align}
        \hat{\psi}(\bm{r}) = \sum_n \phi_n(\bm{r}) \hat{c}_n
    \end{align}
    逆変換は
    \begin{align}
        \hat{c}_n = \int d\bm{r}\ \phi_n^*(\bm{r})\hat{\psi}(\bm{r})
    \end{align}
    で与えられる。この新たな演算子を固有状態の生成・消滅演算子という
    \begin{formula}[固有状態の生成・消滅演算子の交換関係]
        \begin{equation}
            \begin{aligned}
                \{ \hat{c}_n,\hat{c}_m \} &= \{\hat{c}^\dagger_n,\hat{c}_m^\dagger\} = 0 \\
                \{\hat{c}_n,\hat{c}_m^\dagger\} &= \delta_{nm}
            \end{aligned}
        \end{equation}
    \end{formula}
    \begin{derivationBox}
        \begin{align}
            \{ \hat{c}_n,\hat{c}_m \} &= \left\{ \int d\bm{r}\ \phi_n^*(\bm{r})\hat{\psi}(\bm{r}),\int d\bm{r}'\ \phi_m^*(\bm{r}')\hat{\psi}(\bm{r}')  \right\} \notag \\
            &= \int d\bm{r}\int d\bm{r}'\phi_n^*(\bm{r}) \phi_m^*(\bm{r}')\left\{ \hat{\psi}(\bm{r}),\hat{\psi}(\bm{r}')  \right\} \notag \\
            &=0
        \end{align}
        \begin{align}
            \{ \hat{c}^\dagger_n,\hat{c}^\dagger_m \} &= \left\{ \int d\bm{r}\ \phi_n(\bm{r})\hat{\psi}^\dagger(\bm{r}),\int d\bm{r}'\ \phi_m(\bm{r}')\hat{\psi}^\dagger(\bm{r}')  \right\} \notag \\
            &= \int d\bm{r}\int d\bm{r}'\phi_n(\bm{r}) \phi_m(\bm{r}')\left\{ \hat{\psi}^\dagger(\bm{r}),\hat{\psi}^\dagger(\bm{r}')  \right\} \notag \\
            &=0
        \end{align}
        \begin{align}
            \{ \hat{c}_n,\hat{c}^\dagger_m \} &= \left\{ \int d\bm{r}\ \phi_n^*(\bm{r})\hat{\psi}(\bm{r}),\int d\bm{r}'\ \phi_m(\bm{r}')\hat{\psi}^\dagger(\bm{r}')  \right\} \notag \\
            &= \int d\bm{r}\int d\bm{r}'\phi_n^*(\bm{r}) \phi_m(\bm{r}')\left\{ \hat{\psi}(\bm{r}),\hat{\psi}^\dagger(\bm{r}')  \right\} \notag \\
            &=\int d\bm{r}\int d\bm{r}'\phi_n^*(\bm{r}) \phi_m(\bm{r}')\delta(\bm{r}-\bm{r}') \notag \\
            &=\int d\bm{r}\phi_n^*(\bm{r}) \phi_m(\bm{r}) \notag \\
            &= \delta_{nm}
        \end{align}
    \end{derivationBox}
    次にこの演算子を用いてハミルトニアンを書き換える
    \begin{formula}[自由粒子ハミルトニアン]
        \begin{equation}
            \hat{\mathscr{H}}_0 = \sum_n \varepsilon_n\hat{c}^\dagger_n\hat{c}_n
        \end{equation}
    \end{formula}
    \begin{derivationBox}
        \begin{align}
            \hat{\mathscr{H}}_0 &= \int d\bm{r}\left\{ \frac{\hbar^2}{2m} \bm{\nabla}\hat{\psi}^\dagger(\bm{r})\cdot \bm{\nabla}\hat{\psi}(\bm{r}) + V(\bm{r}) \hat{\psi}^\dagger(\bm{r})\hat{\psi}(\bm{r}) \right\} \notag \\
            &= \sum_{n,m}\int d\bm{r}\left\{ \frac{\hbar^2}{2m} \bm{\nabla}\phi_n^*(\bm{r}) \hat{c}^\dagger_n\cdot \bm{\nabla}\phi_m(\bm{r})\hat{c}_m + V(\bm{r}) \phi_n^*(\bm{r}) \hat{c}^\dagger_n\phi_m(\bm{r}) \hat{c}_m \right\} \notag \\
            &= \sum_{n,m}\int d\bm{r}\left\{ -\frac{\hbar^2}{2m} \phi_n^*(\bm{r})  \bm{\nabla}^2\phi_m(\bm{r}) + V(\bm{r}) \phi_n^*(\bm{r}) \phi_m(\bm{r}) \right\} \hat{c}^\dagger_n \hat{c}_m  \notag \\
            &= \sum_{n,m}\int d\bm{r}\ \phi_n^*(\bm{r})\left\{ -\frac{\hbar^2}{2m} \bm{\nabla}^2 + V(\bm{r})  \right\} \phi_m(\bm{r}) \hat{c}^\dagger_n \hat{c}_m  \notag \\
            &= \sum_{n,m}\int d\bm{r}\ \phi_n^*(\bm{r})\varepsilon_m \phi_m(\bm{r}) \hat{c}^\dagger_n \hat{c}_m \notag \\
            &= \sum_{n,m}\delta_{nm} \varepsilon_m  \hat{c}^\dagger_n \hat{c}_m \notag \\
            &= \sum_{n} \varepsilon_n  \hat{c}^\dagger_n \hat{c}_n 
        \end{align}
    \end{derivationBox}
    系の全粒子数演算子は、密度演算子を用いて以下の様に表される。
    \begin{formula}[密度演算子]
        \begin{equation}
            \hat{\mathscr{N}} := \int d\bm{r}\ \hat{\rho}(\bm{r}) = \sum_n \hat{c}_n^\dagger\hat{c}_n
        \end{equation}
    \end{formula}
    \begin{derivationBox}
        \begin{align}
            \hat{\mathscr{N}} &= \int d\bm{r}\ \hat{\rho}(\bm{r}) \notag \\
            &= \int d\bm{r}\ \hat{\psi}^\dagger(\bm{r})\hat{\psi}(\bm{r}) \notag \\
            &= \sum_{n,m}\int d\bm{r}\ \phi_n^*(\bm{r})\phi_m(\bm{r})\hat{c}^\dagger_n\hat{c}_m \notag \\
            &= \sum_{n,m}\delta_{nm} \hat{c}^\dagger_n\hat{c}_m\notag \\
            &=  \sum_n \hat{c}_n^\dagger\hat{c}_n
        \end{align}
    \end{derivationBox}
    一般に$N$粒子状態は
    \begin{align}
        \ket{\Psi_N(T)} = e^{\frac{E}{i\hbar}t} \prod_{i=1}^N \hat{c}^\dagger_{n_i} \ket{0},\quad E = \sum_{i=1}^N \varepsilon_n
    \end{align}
    と書ける。
    \begin{derivationBox}
        
    \end{derivationBox}
\subsection{ハイゼンベルグ描像}
    波動関数の状態ベクトルの時間発展は
    \begin{align}
        i\hbar \frac{\partial}{\partial t}\ket{\Psi_N(t)} = \hat{\mathscr{H}}\ket{\Psi_N(t)}
    \end{align}
    で記述される。これをSchrödinger描像という。\\
    ここでハミルトニアン$\hat{\mathscr{H}}$が時間に依存しないとき
    \begin{align}
        &i\hbar \frac{\partial}{\partial t}\ket{\Psi_N(t)} = \hat{\mathscr{H}}\ket{\Psi_N(t)} \notag \\
        &\Leftrightarrow \ket{\Psi_N(t)} = e^{\frac{\hat{\mathscr{H}}}{i\hbar}t} \ket{\Psi_N(0)} = \hat{U}(t) \ket{\Psi_N(0)}
    \end{align}
    と書くことができ
    \begin{align}
        \hat{U}(t) := e^{\frac{\hat{\mathscr{H}}}{i\hbar}t}
    \end{align}
    と定義し、これを時間発展演算子という。これを用いると物理量$\hat{\mathcal{O}}$の行列要素は
    \begin{align}
        \bra{\Psi'_N(t)} \hat{\mathcal{O}} \ket{\Psi_N(t)} = \bra{\Psi'_N(0)} e^{\frac{-\hat{\mathscr{H}}}{i\hbar}t} \hat{\mathcal{O}} e^{\frac{\hat{\mathscr{H}}}{i\hbar}t}\ket{\Psi_N(0)} 
    \end{align}
    となり、新たに
    \begin{align}
        \hat{\mathcal{O}}_H(t) := e^{-\frac{\hat{\mathscr{H}}}{i\hbar}t} \hat{\mathcal{O}} e^{\frac{\hat{\mathscr{H}}}{i\hbar}t}
    \end{align}
    と定義すると、演算子$\hat{\mathcal{O}}$の行列要素は
    \begin{align}
        \bra{\Psi'_N(0)} \hat{\mathcal{O}}_H(t)\ket{\Psi_N(0)}
    \end{align}
    と表される。このような時間発展を演算子に押し付けた表現をハイゼンベルグ描像という。その時間発展は以下の方程式に従う
    \begin{formula}[ハイゼンベルグ描像の演算子の時間発展]
        \begin{equation}
            i\hbar \frac{\partial}{\partial t} \hat{\mathcal{O}}_H(t) = [\hat{\mathcal{O}}_H(t),\hat{\mathscr{H}}]
        \end{equation}
    \end{formula}
    \begin{derivationBox}
        \begin{align}
            i\hbar \frac{\partial}{\partial t} \hat{\mathcal{O}}_H(t) &= i\hbar \frac{\partial}{\partial t} e^{-\frac{\hat{\mathscr{H}}}{i\hbar}t} \hat{\mathcal{O}} e^{\frac{\hat{\mathscr{H}}}{i\hbar}t} \notag \\
            &= -\hat{\mathscr{H}} e^{-\frac{\hat{\mathscr{H}}}{i\hbar}t} \hat{\mathcal{O}} e^{\frac{\hat{\mathscr{H}}}{i\hbar}t} + e^{-\frac{\hat{\mathscr{H}}}{i\hbar}t} \hat{\mathcal{O}}e^{\frac{\hat{\mathscr{H}}}{i\hbar}t} \hat{\mathscr{H}}  \notag \\
            &= -\hat{\mathscr{H}}  \hat{\mathcal{O}}_H(t)  +  \hat{\mathcal{O}}_H(t) \hat{\mathscr{H}}  \notag \\
            &= [\hat{\mathcal{O}}_H(t),\hat{\mathscr{H}}]
        \end{align}
    \end{derivationBox}
\subsection{電流演算子}
    この節では第二量子化における電流演算子の表式を考える。\\
    電荷保存則
    \begin{align}
        q\frac{\partial}{\partial t} \rho(\bm{r},t) = -\bm{\nabla} \cdot \bm{j}(\bm{r},t)
    \end{align}
    を参考にして、第二量子化における電流密度演算子$\hat{\bm{j}}(\bm{r})$が以下の方程式を満たすこととする
    \begin{align}
        q\frac{\partial}{\partial t} \bra{\Psi_N(t)}\hat{\rho}(\bm{r})\ket{\Psi_N(t)} = -\bm{\nabla} \cdot \bra{\Psi_N(t)} \hat{\bm{j}}(\bm{r}) \ket{\Psi_N(t)}
    \end{align}
    となる。ハイゼンベルグ描像で書き直すと
    \begin{align}
        q\ \frac{\partial}{\partial t} e^{-\frac{\hat{\mathscr{H}}}{i\hbar}t} \hat{\rho}(\bm{r}) e^{\frac{\hat{\mathscr{H}}}{i\hbar}t} = q\frac{\partial}{\partial t} \hat{\rho}_H(\bm{r},t)= -\bm{\nabla}\cdot \hat{\bm{j}}_H (\bm{r},t)
    \end{align}
    となる。さらにハイゼンベルグ描像の演算子の時間発展方程式より
    \begin{align}
        &q\frac{\partial}{\partial t} \hat{\rho}_H(\bm{r},t)= -\bm{\nabla}\cdot \hat{\bm{j}}_H (\bm{r},t) \notag \\
        &\Leftrightarrow \frac{q}{i\hbar} [\hat{\rho}_H(\bm{r},t),\hat{\mathscr{H}}] =  -\bm{\nabla}\cdot \hat{\bm{j}}_H (\bm{r},t) \notag \\
        &\Leftrightarrow \frac{q}{i\hbar} [\hat{\rho}(\bm{r},t),\hat{\mathscr{H}}] =  -\bm{\nabla}\cdot \hat{\bm{j}} (\bm{r},t) 
    \end{align}
    を満たす$\hat{\bm{j}}(\bm{r},t)$が電流密度演算子となる。
\end{document}